\subsection{Influence AA through the Curie--Weiss critical point}\label{ht:full-aa-through-the-curieweiss-critical-point}

\textbf{Theorem HT5.} There is an absolute constant \(c>0\) such that for every \(n\ge 1\), \(0\le \beta\le 1\), and admissible \(f\) of degree at most \(d\ge 1\),

\[
\boxed{\qquad
\max_i I_i^{\mu_{n,\beta}}(f)
\ge c\,\frac{\operatorname{Var}_{\mu_{n,\beta}}(f)^4}{d^{12}}.
\qquad}
\tag{HT5}
\]

In particular, the critical inverse temperature \(\beta=1\) is on the side where distributional AA holds.

The variance exponent four is sharp at this endpoint: the degree-one magnetization \(M_n=(X_1+\cdots +X_n)/n\) has \(I_i(M_n)=n^{-2}\) and \(\operatorname{Var}_{\mu_{n,1}}(M_n)=\Theta(n^{-1/2})\), as proved in Section~\ref{proof:curie_weiss}. A uniform lower bound with variance exponent strictly less than four would already fail on this family. No optimality claim is made for the degree exponent twelve.

\textbf{Latent-measure inputs.} The elementary Hubbard--Stratonovich calculation in Section~\ref{proof:curie_weiss} expresses (CW), for \(0<\beta\le 1\), as the mixture \(\mu=\mathbb E[P_T]\) with \(T=\tanh Z\) and

\[
\nu_{n,\beta}(dz)
\propto \exp\left[-n\left(\frac{z^2}{2\beta}-\log\cosh z\right)\right]dz.
\]

The estimates proved there give absolute constants \(C_0,C_1,\eta>0\) such that:

\begin{enumerate}
\def\labelenumi{\arabic{enumi}.}
\tightlist
\item
  \(T\) is symmetric and \(\mathbb E T^2\le C_0 n^{-1/2}\).
\item
  For some \(a \in (0,1)\), the density of \(T\) is at least \(\eta/a\) on \([-a,a]\).
\item
  Every real function on the cube satisfies \(\operatorname{Var}_\mu f \le  C_1 \sqrt{n} \sum_i I_i^\mu(f)\).
\end{enumerate}

These estimates hold uniformly over \(0<\beta\le 1\); the relevant latent width shrinks as \(\beta\) decreases. They concern the normalized complete-graph model (CW), not all Ising interactions.

\textbf{Proof of HT5.} Let \(v=\operatorname{Var}_\mu f\). If \(v=0\), the conclusion is immediate. Assume \(v>0\), and set \(B=18C_0\).

If \(n \ge  B^2 d^8/v^2\), equation (4.1) implies \(\operatorname{Var}_U f\ge v/2\). By HT4 and (UAA),

\[
\max_i I_i^\mu(f)
\ge c\frac{v^{5/2}}{d^8}
\ge c\frac{v^4}{d^{12}},
\]

where the last inequality uses \(0<v\le 1\) and \(d\ge 1\).

If \(n < B^2 d^8/v^2\), the latent Poincaré bound gives

\[
\max_i I_i^\mu(f)
\ge\frac{v}{C_1 n^{3/2}}
>\frac{v^4}{C_1B^3d^{12}}.
\]

This proves the uniform claim for \(0<\beta\le 1\). At \(\beta=0\), (UAA) directly implies HT5. \(\square\)
